% Problem record op_b54ea1af90e24eaa. This TeX file is derived from
% database/problems_json/op_b54ea1af90e24eaa.json by scripts/sync-tex.mjs; edit the JSON record
% and rerun the script rather than editing this file.
\section{Semialgebraicity of closed quantum correlations}
\label{sec:op_b54ea1af90e24eaa}

\paragraph{Problem.}
Is the closure of the finite-dimensional tensor-product correlation set
semialgebraic in every fixed bipartite Bell scenario?  Fix positive integers
$n_A,n_B,m_A,m_B$.  Inputs satisfy $1\leq x\leq n_A$ and
$1\leq y\leq n_B$, while outputs satisfy $1\leq a\leq m_A$ and
$1\leq b\leq m_B$.  The set
$\mathcal C_q(n_A,n_B,m_A,m_B)$ consists of behaviors
\begin{equation}
  p(a,b\mid x,y)
  =\operatorname{Tr}\!\left[
    \rho\bigl(E_a^x\otimes F_b^y\bigr)
    \right]
  \label{eq:p31-quantum-behavior}
\end{equation}
obtained from arbitrary finite-dimensional local Hilbert spaces, a density
operator $\rho$, and local POVMs satisfying
\begin{equation}
  E_a^x\succeq0,\qquad F_b^y\succeq0,\qquad
  \sum_{a=1}^{m_A}E_a^x=I,\qquad
  \sum_{b=1}^{m_B}F_b^y=I.
  \label{eq:p31-povm-constraints}
\end{equation}
Equation~\eqref{eq:p31-povm-constraints} fixes the measurement class in
Eq.~\eqref{eq:p31-quantum-behavior}.  Define the closed correlation set by
\begin{equation}
  \mathcal C_{qa}(n_A,n_B,m_A,m_B)
  :=\overline{\mathcal C_q(n_A,n_B,m_A,m_B)}.
  \label{eq:p31-closed-correlation-set}
\end{equation}
For every fixed tuple, is the set in
Eq.~\eqref{eq:p31-closed-correlation-set} a finite Boolean combination of
polynomial equalities and inequalities with real coefficients?  The
coefficients may be arbitrary real numbers, and the description need not be
one conjunction of weak inequalities.

\subsection*{Status}

Unsolved

\subsection*{Source}

Tsirelson explicitly asks whether the quantum-correlation set in each fixed
finite Bell scenario is semialgebraic
\sourcecite{ref:p31-tsirelson}{Tsi93}.

\subsection*{Progress}

\begin{itemize}
  \item In the scenario with two binary measurements per party, Jordan-type
  dimension reduction and finite convexification give a semialgebraic
  parametrization \sourcecite{ref:p31-tsirelson}{Tsi93},
  \sourcecite{ref:p31-masanes}{Mas05}.  This does not extend to every fixed
  input--output tuple.

  \item The zero-marginal correlator slice in the smallest bipartite scenario
  has an explicit semialgebraic description but is not basic semialgebraic.
  Consequently, asking for one finite conjunction of polynomial inequalities
  would already be too strong \sourcecite{ref:p31-le}{LMSWZ23}.

  \item The finite-dimensional set $\mathcal C_q$ defined by
  Eqs.~\eqref{eq:p31-quantum-behavior}--\eqref{eq:p31-povm-constraints} is not
  closed in sufficiently large fixed scenarios.  This necessitates the
  closure in Eq.~\eqref{eq:p31-closed-correlation-set}
  \sourcecite{ref:p31-slofstra}{Slo19}.

  \item For sufficiently large fixed input and output sizes, membership in
  $\mathcal C_{qa}$ is undecidable even for behaviors with algebraic
  coordinates.  This excludes effective polynomial descriptions with
  computable algebraic coefficients, but it does not exclude a non-effective
  semialgebraic presentation using arbitrary real coefficients
  \sourcecite{ref:p31-fu}{FMS25}.
\end{itemize}

\subsection*{References}

\begin{enumerate}[label={},leftmargin=4.8em,labelsep=0.5em]
  \footnotesize
  \item[\textup{[Tsi93]}]\label{ref:p31-tsirelson}
  B. S. Tsirelson, ``Some Results and Problems on Quantum Bell-Type
  Inequalities,'' \emph{Hadronic Journal Supplement} \textbf{8}(4), 329--345
  (1993).
  \href{https://cris.tau.ac.il/en/publications/some-results-and-problems-on-quantum-bell-type-inequalities/}{publication record};
  \href{https://ma.huji.ac.il/~ohadfeld/Tsirelson/download/hadron.html}{full text}.

  \item[\textup{[Mas05]}]\label{ref:p31-masanes}
  Ll. Masanes, ``Extremal Quantum Correlations for $N$ Parties with Two
  Dichotomic Observables per Site,'' arXiv:quant-ph/0512100 (2005).
  \href{https://doi.org/10.48550/arXiv.quant-ph/0512100}{doi:10.48550/arXiv.quant-ph/0512100};
  \href{https://arxiv.org/abs/quant-ph/0512100}{arXiv:quant-ph/0512100}.

  \item[\textup{[LMSWZ23]}]\label{ref:p31-le}
  T. P. Le, C. Meroni, B. Sturmfels, R. F. Werner, and T. Ziegler, ``Quantum
  Correlations in the Minimal Scenario,'' \emph{Quantum} \textbf{7}, 947
  (2023). \href{https://doi.org/10.22331/q-2023-03-16-947}{doi:10.22331/q-2023-03-16-947};
  \href{https://arxiv.org/abs/2111.06270}{arXiv:2111.06270}.

  \item[\textup{[Slo19]}]\label{ref:p31-slofstra}
  W. Slofstra, ``The Set of Quantum Correlations Is Not Closed,''
  \emph{Forum of Mathematics, Pi} \textbf{7}, e1 (2019).
  \href{https://doi.org/10.1017/fmp.2018.3}{doi:10.1017/fmp.2018.3};
  \href{https://arxiv.org/abs/1703.08618}{arXiv:1703.08618}.

  \item[\textup{[FMS25]}]\label{ref:p31-fu}
  H. Fu, C. A. Miller, and W. Slofstra, ``The Membership Problem for
  Constant-Sized Quantum Correlations Is Undecidable,'' \emph{Communications
  in Mathematical Physics} \textbf{406}, 96 (2025).
  \href{https://doi.org/10.1007/s00220-024-05229-7}{doi:10.1007/s00220-024-05229-7};
  \href{https://arxiv.org/abs/2101.11087}{arXiv:2101.11087}.
\end{enumerate}

\subsection*{Comment}

Undecidable membership rules out effective descriptions of the usual kind,
but semialgebraicity with unrestricted real coefficients is a set-theoretic
existence question and remains unresolved for general fixed scenarios.

\subsection*{Field}

Quantum foundations; Mathematics of quantum information

\subsection*{Topic}

Bell nonlocality; Convex geometry; Operator algebras

\subsection*{ID}

\texttt{op\_b54ea1af90e24eaa}
